%!TEX root=../../main.tex

Let \( X \in \R^{n \times d_0} \) be a matrix of observations and \( y \in
\R^{n \times c} \) a matrix of \( c \)-dimensional responses for each
observation.
This setup covers both multi-class classification and regression with
vector-valued targets.
We are interested in fully-connected ReLU neural networks with arbitrarily many
layers.
Let \( \Wl \in \R^{d_{l-1} \times d_{l}} \) be the matrix of weights for the \(
^{\text{th}} \) layer.
We will use \( \Wl_{i} \) to refer to the \( i^{\text{th}} \) row of \( \Wl \)
and \( \Wl_{\cdot, j} \) to refer to the \( j^{\text{th}} \) column.
The activations at layer $l$ are defined to be,
\begin{equation}\label{eq:activations}
    Z^{(l+1)} = \rbr{\Z W^{(l)}}_+,
\end{equation}
where we take \( Z^{(1)} = X \) and \( \rbr{\cdot}_+ \) is ReLU activation
function defined element-wise as \( \rbr{z}_+ = \max\cbr{z, 0} \).
Using this notation, a fully-connected neural network with \( L \)-layers has
the prediction function given by,
\begin{equation}\label{eq:predictions}
    f_\theta(X) = Z^{(L)} W^{(L)},
\end{equation}
where \( W^{(L)} \in \R^{d_{L-1} \times c} \) ensures
that the prediction is of the right dimension.
Here, \( \theta = \cbr{W^{(1)}, \ldots, \Wk} \) denotes the collection of
weights for the network and \( \Theta \) the complete parameter space.
The choice of \( L = 2 \) corresponds to a two-layer ReLU network,
which is sometimes referred to as a network with one hidden layer.
In this case, the recursive definition can be unrolled to give a simple
and explicit prediction function,
\begin{equation}\label{eq:two-layer-vector-output}
    f_\theta(X) = \rbr{X W^{(1)}}_+ W^{(2)}.
\end{equation}
For simplicity, when \( L = 2 \) we will often use \( d \) for the input
dimension (instead of \( d_0 \)) and \( m \) as the number of hidden units
(instead of \( d_1 \)).

Many of our results are focused on two-layer ReLU networks with scalar
outputs, meaning \( c = 1 \).
In this setting, it will often be useful to break-down \( f_\theta \) into
an explicit sum over neurons in the hidden layer,
\begin{equation}\label{eq:two-layer-scalar-output}
    f_\theta(X) = \sum_{j = 1}^{m} \rbr{X W^{(1)}_{\cdot, j}}_+ W^{(2)}_{j}.
\end{equation}
Note that since \( c = 1 \), there is no weight sharing and
\( W^{(2)} \in \R^{m \times 1} \) is a column vector.

We consider choosing \( \theta \) by solving a regularized empirical risk
minimization problem using the fixed training set \( (X, y) \).
That is, problems of the form,
\begin{equation}\label{eq:regularized-erm}
    \theta^* \in \argmin_{\theta \in \Theta}
    \calL\rbr{f_\theta(X), y} +
    \lambda R(\theta),
\end{equation}
where \( \calL \) is a strictly convex loss function, \( R \) is a (typically
convex) regularizer, and \( \lambda \geq 0 \) controls the strength of the
regularization.
The function
\begin{equation}
    F_\lambda(\theta) = \calL(f_\theta(X), y) + \lambda R(\theta)
\end{equation}
is referred to as the objective function.
Although \( \calL \) and \( R \) are usually convex, \( \theta \mapsto
F_\lambda(\theta) \) is not because the neural network \( f_\theta \) is
highly non-linear and itself non-convex.
As a result, Problem~\ref{eq:regularized-erm} is challenging to solve in both
theory and practice.

The most common choice of \( \calL \) and \( \R \) in practice are squared
error and \( \ell_2 \)-squared (weight-decay) regularization, respectively.
This leads to the following concrete optimization problem:
\begin{equation}\label{eq:squared-error}
    \theta^* \in \argmin_{\theta \in \Theta}
    \half \norm{f_\theta(X) - y}_2^2 +
    \frac{\lambda}{2} \sum_{l=1}^L \norm{\Wl}_F^2,
\end{equation}
where \( \norm{\cdot}_F \) is the Frobenius norm given by the sum of squared entries,
\[
    \norm{W}_F = \rbr{\sum_{i, j} W_{i, j}^2}^{1/2}.
\]
The squared-error loss function is equivalent to assuming a Gaussian likelihood
function and is a standard choice for regression problems.
Recent work has also shown that it performs comparably to cross-entropy for
classification problems \citep{hui2021evaluation}.
Similarly, \( \ell_2 \)-squared regularization is comparable to assuming
a Gaussian prior on the model weights;
from an optimization perspective, it is desirable because it gives curvature to
objective function, ensuring that \( F_\lambda \) is coercive and that the
level sets are compact.

\begin{figure}[t]
    \centering
    %! TEX root = ../../main.tex

%% Illustration of cone decomposition. 

\begin{tikzpicture}[scale=1,
        declare function={
                cone_1(\x)= 3*\x;
                cone_2(\x)= -\x;
                cone_3(\x)= -4*\x;
                bounds(\x)= \x - 10;
            }
    ]
    \begin{axis}[width=\linewidth, height=8cm,
            axis lines=center, yticklabels={,,}, xticklabels={,,},
            ymin=-4, ymax=4, ytick={-4,...,4}, ylabel=$$, x axis line style={-},
                xmin=-6, xmax=6, xtick={-6,...,6}, xlabel=$$, y axis line style={-},
        ]
        \draw [->, draw=red, line width = 0.8mm] (axis cs:0,0) -- (axis cs:3,-2) node[pos=0.6, above] {$W^{(1)}_{\cdot, j}$};
        \addplot[name path=cone_1, domain=-6:6, samples=100, line width=1pt]{cone_1(x)};
        \addplot[name path=bounds, domain=-6:6, samples=100, line width=1pt]{bounds(x)};

        % add color fill to both cones.

        \addplot fill between[
                of = cone_1 and bounds,
                %split, % calculate segment
                every even segment/.style = {fill=green, fill opacity=0.3},
            ];

        %% point labels

        % active examples 
        \node[label=right:$x_1$, circle, fill, inner sep=0.75mm] at (axis cs:2,2) {};
        \node[label=right:$x_4$, circle, fill, inner sep=0.75mm] at (axis cs:1.5,-2) {};

        % inactive examples 
        \node[label=right:$x_2$, circle, line width=0.4mm, draw=black, inner sep=0.75mm] at (axis cs:-3,-2) {};
        \node[label=right:$x_3$, circle, line width=0.4mm, draw=black, inner sep=0.75mm] at (axis cs:-4,-1) {};
        \node[label=right:$x_5$, circle, line width=0.4mm, draw=black, inner sep=0.75mm] at (axis cs:0.5,3.5) {};

        % activation pattern
        \node[] at (axis cs:-3.25,2) {
            $D_j = \begin{bmatrix}
                    1 & 0 & 0 & 0 & 0 \\
                    0 & 0 & 0 & 0 & 0 \\
                    0 & 0 & 0 & 0 & 0 \\
                    0 & 0 & 0 & 1 & 0 \\
                    0 & 0 & 0 & 0 & 0 \\
                \end{bmatrix}
            $};

        % lines
        %\draw [->, draw=red, line width = 0.4mm] (axis cs:2,-2) -- (axis cs:3,-1) node[midway,below right] {$S_{22} \cdot x_2$};
        %\draw [->, draw=red, line width = 0.4mm] (axis cs:0.5,-2) -- (axis cs:-2.5,-2.75) node[pos=0.9,below right] {$S_{33} \cdot x_3$};

        % origin point
        %\node[circle, fill, inner sep=2pt] at (axis cs:0,0) {};

    \end{axis}

\end{tikzpicture}%

    \caption{Illustration of the activation pattern \( D_j \) associated with a
        ReLU neuron parameterized by \( W^{(1)}_{\cdot, j} \).
        The neuron is associated with a half-space where it
        is active (light green) and a half-space where it is inactive.
        The activation pattern \( D_j \) is a diagonal matrix encoding the
        active (solid circles) and inactive (hollow circles) data points:
        \( \sbr{D_{j}}_{ii} = 1 \) if \( x_i \) is active and
        \( \sbr{D_{j}}_{ii} = 0 \) otherwise.
        %This linearizes the ReLU activation on \( \calK_j \): for every
        %\( W^{(1)}_{\cdot, l} \in \calK_j \),
        %\( D_{j} X W^{(1)}_{\cdot, l} = \rbr{X W^{(1)}_{\cdot, l}}_+ \).
    }%
    \label{fig:activation-pattern}
\end{figure}

We conclude by introducing background on convex reformulations.
In the following, we assume that \( \theta \) parameterizers a two-layer,
scalar output ReLU network as described in \cref{eq:two-layer-scalar-output}.
For such networks, convex reformulations re-write \cref{eq:regularized-erm} as
a convex program by enumerating the activations a single neuron in the
hidden layer can take on for fixed \( X \) as follows:
\begin{equation}\label{eq:activation-patterns}
    \calD_X = \cbr{D = \diag(\mathds{1}(X u \geq 0)) : u \in \R^d}.
\end{equation}
This set grows as \(|\calD_X| \le O(r(n/r)^r)\), where \(r := \text{rank}(X)\)
\citep{pilanci2020convexnn}.
Each ``activation pattern'' \( D_i \in \calD_{X} \)
(see \cref{fig:activation-pattern} for a visualization) is associated with
a convex cone,
\begin{equation}\label{eq:activation-cone}
    \calK_i = \cbr{u \in \R^d : (2D_i - I) X u \succeq 0}.
\end{equation}
If \( u \in \calK_i \), then \( u \) matches \( D_i \), meaning \( D_i X u = \rbr{X u}_+ \).
For any subset \( \tilde \calD \subseteq \calD_X \),
the convex reformulation is,
\begin{equation}\label{eq:convex-relu-mlp}
    \begin{aligned}
        \min_{v, w} \, & \calL\Big(\sum_{D_i \in \tilde \calD} D_i X (v_i - u_i), y\Big)+\lambda\sum_{D_i \in \tilde \calD}\norm{v_i}_2+\norm{u_i}_2 \\
                       & \text{s.t.} \quad v_i, u_i \in \calK_i.
    \end{aligned}
\end{equation}
\begin{figure}[t]
    \centering
    \begin{tikzpicture}[scale=1,
    ]
    \begin{axis}[width=0.6\linewidth, height=5.0cm,
            axis lines=none,  % don't print axis lines
            yticklabels={,,}, xticklabels={,,},
            ymin=0, ymax=8, y axis line style={-},
            xmin=0, xmax=15, x axis line style={-},
        ]
        \node [] (input) at (axis cs:1,4) {$X$};
        \node [Input] (input1) at (axis cs:3,2) {};
        \node [Input] (input2) at (axis cs:3,4) {};
        \node [Input] (input3) at (axis cs:3,6) {};

        \draw [decorate, line width = 0.6mm,
            decoration = {calligraphic brace}] (axis cs:2,1.5) --  (axis cs:2,6.5);

        \node [Hidden] (hidden1) at (axis cs:8,1) {};
        \node [Hidden] (hidden2) at (axis cs:8,3) {};
        \node [Hidden] (hidden3) at (axis cs:8,5) {};
        \node [Hidden] (hidden4) at (axis cs:8,7) {};

        \draw [->, style=arrow, draw=black] (input1) -- (hidden1);
        \draw [->, style=arrow, draw=black] (input2) -- (hidden1);
        \draw [->, style=arrow, draw=black] (input3) -- (hidden1);

        \draw [->, style=arrow, draw=black] (input1) -- (hidden2);
        \draw [->, style=arrow, draw=black] (input2) -- (hidden2);
        \draw [->, style=arrow, draw=black] (input3) -- (hidden2);

        \draw [->, style=arrow, draw=black] (input1) -- (hidden3);
        \draw [->, style=arrow, draw=black] (input2) -- (hidden3);
        \draw [->, style=arrow, draw=black] (input3) -- (hidden3);

        \draw [->, style=arrow, draw=red] (input1) -- (hidden4);
        \draw [->, style=arrow, draw=red] (input2) -- (hidden4);
        \draw [->, style=arrow, draw=red] (input3) -- (hidden4) node[pos=0.5,above] {$W_{1j}$};

        \node [Output] (output) at (axis cs:13,4) {$\hat y$};

        \draw [->, style=arrow, draw=black] (hidden1) -- (output);
        \draw [->, style=arrow, draw=black] (hidden2) -- (output);
        \draw [->, style=arrow, draw=black] (hidden3) -- (output);
        \draw [->, style=arrow, draw=red] (hidden4) -- (output) node[pos=0.5,above] {$w_{2j}$};
    \end{axis}

\end{tikzpicture}%
\begin{tikzpicture}[scale=1,
    ]
    \begin{axis}[width=0.6\linewidth, height=6.5cm,
            axis lines=none,  % don't print axis lines
            yticklabels={,,}, xticklabels={,,},
            ymin=0, ymax=8, y axis line style={-},
            xmin=2, xmax=14, x axis line style={-},
        ]
        \node [] (input1) at (axis cs:3,1) {$X$};
        \node [] (input2) at (axis cs:3,2.75) {$X$};
        \node [] (input3) at (axis cs:3,4.5) {$X$};

        \node [Splits] (hidden1) at (axis cs:9,1) {\scriptsize $D_3 X$};
        \node [Splits] (hidden2) at (axis cs:9,2.75) {\scriptsize $D_2 X$};
        \node [Splits] (hidden3) at (axis cs:9,4.5) {\scriptsize $D_1 X$};

        \draw [draw=blue, dashed, line width=0.4mm] (axis cs:6.55,6.05) -- (axis cs:8.1,6.85);

        \draw [->, style=arrow, draw=black] (input1) -- (hidden1);

        \draw [->, style=arrow, draw=black] (input2) -- (hidden2);

        \draw [->, style=arrow, draw=red] (input3) -- (hidden3);

        \node [Output] (output) at (axis cs:13,2.75) {$\hat y$};

        \draw [->, style=arrow, draw=black] (hidden1) -- (output);
        \draw [->, style=arrow, draw=black] (hidden2) -- (output);
        \draw [->, style=arrow, draw=red] (hidden3) -- (output) node[pos=0.5,above] {$u_j$};

        \draw [draw=black, line width=0.3mm, dashed] (axis cs:5.7, 4.235) -- (axis cs:7.75,5.755);
        \draw [draw=black, line width=0.3mm, dashed] (axis cs:5.28, 4.75) -- (axis cs:6.55,6.95);

        \node [draw=black, minimum size=0.3cm, shape=circle, solid, line width=0.5mm] (examine) at (axis cs:5.5,4.5) {};
        \node [draw=black, minimum size=1.1cm, shape=circle, solid, line width=0.5mm] (closeup) at (axis cs:7.3,6.5) {};

        \node [fill=white, draw=gray, line width=0.3mm, inner sep=0.05cm, shape=circle] at (axis cs:7.7,7.0) {};
        \node [fill=white, draw=gray, line width=0.3mm, inner sep=0.05cm, shape=circle] at (axis cs:7.15,6.9) {};
        \node [fill=white, draw=gray, line width=0.3mm, inner sep=0.05cm, shape=circle] at (axis cs:7.3,7.2) {};
        \node [fill=white, draw=gray, line width=0.3mm, inner sep=0.05cm, shape=circle] at (axis cs:6.9,6.5) {};

        \node [fill=red, inner sep=0.06cm, shape=circle] at (axis cs:6.9,5.95) {};
        \node [fill=red, inner sep=0.06cm, shape=circle] at (axis cs:7.4,6.0) {};
        \node [fill=red, inner sep=0.06cm, shape=circle] at (axis cs:7.8,6.4) {};

    \end{axis}

\end{tikzpicture}%

    \caption{Comparison between a two-layer ReLU network and its convex
        reformulation.
        The non-convex ReLU network maps the input \( X \) to a hidden
        representation and then predicts using a linear combination of these
        neurons.
        In comparison, the convex reformulation computes \( \hat{y} \) as a sum
        of linear models fit to copies of \( X \) which are split into active
        and inactive data-points by half-spaces passing
        through the origin.}%
    \label{fig:convex-reformulation}
\end{figure}
See \cref{fig:convex-reformulation} for an illustration.
\citet{pilanci2020convexnn} prove that this program and \cref{eq:non-convex-relu-mlp}
are equivalent in the following sense:
if \( \tilde \calD = \calD_X \) and \( m \geq m^* \) for some \( m^* \leq n + 1 \),
then the two programs have the same optimal value and every solution to the convex program
can be mapped to a solution of the non-convex training problem and vice-versa.
Given a solution \( \rbr{v^*, u^*} \), optimal weights
for the ReLU problem are given by
\begin{equation}\label{eq:convex-to-relu}
    \begin{aligned}
        W^{(1)}_{\cdot, i} & = v_i^*/ \sqrt{\norm{v_i^*}}, \quad W^{(2)}_{i} = \sqrt{\norm{v_i^*}}
        \\
        W_{1j}             & = u_i^*/ \sqrt{\norm{u_i^*}}, \quad w_{2j} = -\sqrt{\norm{u_i^*}},
    \end{aligned}
\end{equation}
where we use the convention that \( 0/0 = 0 \).

In practice, learning with \( \calD_{X} \) is intractable except when the data
are low rank.
In the following chapter, we provide refined conditions on \( \tilde \calD \)
which are sufficient for \cref{eq:convex-relu-mlp} to be equivalent to the
non-convex problem.

